\scanpage{56}
ce qui fait qu'il suffit de définir l'un des deux isom. (34) -- et en vertu de la définition de $\Pi_{03}$ en (32), il est naturel de définir plutôt l'iso
\[
\deleted{(\unsure{35})}\qquad\psi_r=\psi_{\deleted{\unsure{\omega}}\varepsilon}^{i,\omega}:\Pi_{03}\xrightarrow{\sim}\Pi_\varepsilon
\]
\margin{$r=(i,\omega,\deleted{\unsure{\varepsilon}}\inserted{\unsure{\bar\omega|\varpi}})$}
par transport des structures comme
\[
(36)\quad\left\{\begin{aligned}
\psi_{\varepsilon\deleted{\unsure{\omega}}}^{i,\omega}&=\deleted{\Pi_1(\unsure{u_\varepsilon^{i,\omega}})}\Pi_1(u_{i,\omega,\varepsilon};P^+)\\
\text{où }\inserted{u_r=u_{i,\omega,\varepsilon}}\ u_\varepsilon^{i,\omega}&:(U_{0,3}(\mathbb C),P_+)\longrightarrow(\Sigma^*,P_\omega)
\end{aligned}\right.
\]
est l'unique iso conforme qui transforme le repère \deleted{\unsure{base}} standart \inserted{$(\infty,\omega_{0,1,\infty},+1)$} \deleted{\unsure{analogue}} dans le repère \deleted{$\unsure{(\omega i,\omega'i,i)}$} $(i,\omega,\varepsilon)$
\margin{Au couple $(i,\omega)$ on associe le repère (au sens habituel) $r=(\omega i,\omega'i,i)$ de $J$, i.e. \deleted{\unsure{correspondant}}
\[
0\longmapsto\omega i,\quad 1\longmapsto\omega^2i=\omega'i,\quad\infty\longmapsto i.
\]
Ainsi, on a, si $r=\deleted{\unsure{(\omega i,\omega'i,i)}}(i,\omega)$,
\[
\rho_\omega(r)=r\circ\rho=(\omega i,\omega).
\]
}
Pour un $P_\omega$ fixé, il y a six choix pour les \inserted{$\psi_r=$} $\psi_{\varepsilon\deleted{\unsure{\omega}}}^{i,\omega}:\Pi_{03}\longrightarrow\Pi_\varepsilon$ -- et pour un $Q_i$ fixé, il y a quatre choix pour $\omega,\varepsilon$. Cela fait \inserted{trois resp.} deux choix si on suppose que l'on a choisi une \unsure{valeur} de $\unsure{\bar\omega|\varpi}$ sur $\Sigma$ et réalisé \unsure{donc} $\uline{\varepsilon}\simeq\uline{\omega}$, et si on se borne aux $\varepsilon$ qui correspondent à $\omega\ \unsure{\Rightarrow}$ -- il faut donc voir comment passer de l'un \inserted{de ces choix} à l'autre.
\margin{55}
